\section*{Chapter \ref{ch:b2:anthropogenic-climate-change} --- Anthropogenic Climate Change}

\begin{solution}{exo:b2:anthropogenic-climate-change:1}
(a) $3.75 - 0.60 + 0.22 + 0.45 + 0.11 = \qty{3.93}{W.m^{-2}}$, the adjustments adding
\qty{0.18}{W.m^{-2}}. (b) $3.93/3.71 = 1.059$: \qty{5.9}{\percent} more. (c) The tropospheric
temperature adjustment, $-0.60$: a warmer troposphere emits more infrared to space, which removes
part of the extra energy before the surface has responded.
\end{solution}

\begin{solution}{exo:b2:anthropogenic-climate-change:2}
(a) $4\pi(\num{6.371e6})^2 = \qty{5.10e14}{m^2}$. (b) Over 47 years, $\num{434.9e21}/(47 \times
\num{3.156e7} \times \num{5.10e14}) = \qty{0.575}{W.m^{-2}}$; over 12 years, $\num{152.4e21}/(12
\times \num{3.156e7} \times \num{5.10e14}) = \qty{0.789}{W.m^{-2}}$, the 0.57 and 0.79 of the
assessment. (c) $0.91 \times 434.9 = \qty{396}{ZJ}$.
\end{solution}

\begin{solution}{exo:b2:anthropogenic-climate-change:3}
(a) $-3.22 + 1.30 + 0.35 + 0.42 - 0.01 = \num{-1.16}$, the net. (b) $3.93/1.16 =
\qty{3.39}{\celsius}$; $3.93/3.22 = \qty{1.22}{\celsius}$. (c) $3.39/1.22 = 2.78$: the feedbacks
almost triple the Planck warming.
\end{solution}

\begin{solution}{exo:b2:anthropogenic-climate-change:4}
(a) One petagram of carbon is $44.009/12.011 = 3.664$~Gt of carbon dioxide: $0.45 \times 3.664 =
\qty{1.65}{\celsius}$ per \qty{1000}{PgC}, and the range becomes 0.989 to 2.31. (b) $2390 \times
12.011/44.009 = \qty{652}{PgC}$. (c) $0.45 \times 2.39 = \qty{1.08}{\celsius}$.
\end{solution}

\begin{solution}{exo:b2:anthropogenic-climate-change:5}
(a) $g = 159/398 = 0.400$; $T = (398/\num{5.670e-8})^{1/4} = \qty{289}{K}$. (b) $-0.6 \times 4
\times \num{5.670e-8} \times 288^3 = \qty{-3.25}{W.m^{-2}.K^{-1}}$; at \qty{298}{K}, $-3.60$.
(c) Less sensitive: the response grows as $T^3$, so a warmer planet sheds each extra degree more
easily, other feedbacks being equal.
\end{solution}

\begin{solution}{exo:b2:anthropogenic-climate-change:6}
(a) $1.85 - 0.50 = \qty{1.35}{W.m^{-2}.\celsius^{-1}}$. (b) $\sqrt{0.32^2 + 0.20^2} = 0.38$.
(c) The sum is far better known than either term: their errors are strongly anticorrelated. Where
the upper troposphere warms more, it also holds more vapour, so a stronger negative \omterm{def:b2:anthropogenic-climate-change:feedbacks}{lapse-rate
feedback} comes with a stronger positive \omterm{def:b2:anthropogenic-climate-change:feedbacks}{water-vapour feedback}.
\end{solution}

\begin{solution}{exo:b2:anthropogenic-climate-change:7}
(a) $\lambda = (0.59 - 2.20)/1.03 = \qty{-1.56}{W.m^{-2}.\celsius^{-1}}$; $3.93/1.56 =
\qty{2.52}{\celsius}$. (b) $3.93/(1.56 - 0.5) = 3.93/1.06 = \qty{3.71}{\celsius}$. (c) $3.93/1.46 =
\qty{2.69}{\celsius}$. The assessment, with full distributions, finds 3.5 and \qty{2.7}{\celsius}.
\end{solution}

\begin{solution}{exo:b2:anthropogenic-climate-change:8}
(a) $0.617/3682 = \qty{1.68e-4}{\celsius^{-1}}$. (b) $0.113 \times 0.396 = \qty{0.0447}{m}$, or
\qty{44.7}{mm}, close to the \qty{47.5}{mm} of the budget. (c) $\num{2e-4} \times 0.2 \times 700 =
\qty{0.028}{m}$, \qty{28}{mm}.
\end{solution}

\begin{solution}{exo:b2:anthropogenic-climate-change:9}
(a) $\sum xy = 0.311$, $\sum x^2 = 0.300$: $\beta = 1.04$. (b) The \qty{95}{\percent} interval is
$1.04 \pm 1.96 \times 0.08$, 0.88 to 1.19: it excludes 0, so the response is detected, and contains
1, so the observed change is consistent with the simulated one. (c) The interval 0.34 to 0.66 still
excludes 0 (detected), but excludes 1: the model would overestimate the response by about a factor
of two.
\end{solution}

\begin{solution}{exo:b2:anthropogenic-climate-change:10}
(a) $\sum g^2 = 1.20$, $\sum a^2 = 0.18$, $\sum ga = -0.46$, $\sum gy = 0.74$, $\sum ay = -0.28$.
The normal equations are $1.20\,\beta_g - 0.46\,\beta_a = 0.74$ and $-0.46\,\beta_g + 0.18\,\beta_a
= -0.28$; their determinant is $1.20 \times 0.18 - 0.46^2 = 0.0044$, and $\beta_g = (0.74 \times
0.18 - 0.46 \times 0.28)/0.0044 = 1.00$, $\beta_a = (-1.20 \times 0.28 + 0.46 \times 0.74)/0.0044 =
1.00$. (b) Greenhouse gases $+\qty{0.80}{\celsius}$, aerosols $-\qty{0.30}{\celsius}$. (c) The
determinant is small because $g$ and $a$ are nearly proportional (their correlation is $-0.95$): a
little noise in $y$ can be absorbed by raising both scaling factors together, so each is poorly
constrained though their combined effect is well known.
\end{solution}

\begin{solution}{exo:b2:anthropogenic-climate-change:11}
(a) $(0.93 - 0.1)/0.45 \times 1000 = \qty{1844}{Gt}$; minus 210 and $26 \times 0.93 = 24$:
\qty{1610}{Gt}. (b) $0.93 - x = 0.45 \times (1350 + 210 + 24)/1000 = 0.713$, so $x =
\qty{0.22}{\celsius}$: about the upper end of the assessed 0.1 to \qty{0.2}{\celsius}. (c) The
budget depends on the non-CO$_2$ warming $x$ only through $-x/\text{TCRE}$: each
\qty{0.1}{\celsius} costs $0.1/0.45 \times 1000 = \qty{222}{Gt}$, whatever the limit.
\end{solution}

\begin{solution}{exo:b2:anthropogenic-climate-change:12}
(a) $3.93/1.81 = \qty{2.17}{\celsius}$, $3.93/0.51 = \qty{7.71}{\celsius}$, $3.93/1.16 =
\qty{3.39}{\celsius}$. (b) The low side is $3.39 - 2.17 = \qty{1.22}{\celsius}$ long, the high side
$7.71 - 3.39 = \qty{4.32}{\celsius}$: the function $1/|\lambda|$ is convex, so equal steps of
$\lambda$ towards zero raise the ECS more and more. (c) A weakly negative $\lambda$ gives a very
large ECS, and small changes of $\lambda$ near zero are hard to exclude; a low ECS needs a strongly
negative $\lambda$, which the \omterm{def:b2:anthropogenic-climate-change:feedbacks}{Planck response} bounds and the observed warming contradicts.
\end{solution}

\begin{solution}{pb:b2:anthropogenic-climate-change:1}
\textbf{1.} $2.16 + 0.54 + 0.21 + 0.41 = \qty{3.32}{W.m^{-2}}$.
\textbf{2.} $-0.22 - 0.84 = \qty{-1.06}{W.m^{-2}}$.
\textbf{3.} $3.32 + 0.47 + 0.05 - 1.06 - 0.20 + 0.08 + 0.06 = \qty{2.72}{W.m^{-2}}$, the assessed
total.
\textbf{4.} $3.93\log_2(410/278) = \qty{2.20}{W.m^{-2}}$, within \qty{2}{\percent} of 2.16.
\textbf{5.} $3.93\log_2(427.35/278) = \qty{2.44}{W.m^{-2}}$.
\textbf{6.} The aerosol--cloud term, $-1.45$ to $-0.25$, \qty{1.20}{W.m^{-2}} wide, against 1.52 for
the total: it carries most of the uncertainty of the total because it is large and poorly
constrained.
\textbf{7.} $0.79 \times \num{5.10e14} \times 12 \times \num{3.156e7} = \qty{1.53e23}{J}$,
\qty{153}{ZJ}.
\textbf{8.} $0.91 \times 153 = \qty{139}{ZJ}$.
\textbf{9.} $\lambda = (0.59 - 2.20)/1.03 = \qty{-1.56}{W.m^{-2}.\celsius^{-1}}$.
\textbf{10.} $3.93/1.56 = \qty{2.52}{\celsius}$.
\textbf{11.} $3.93 \times 1.03/2.20 = \qty{1.84}{\celsius}$.
\textbf{12.} $3.93/1.16 = \qty{3.39}{\celsius}$, higher by \qty{0.87}{\celsius}.
\textbf{13.} The warming so far has been strongest where the feedbacks are most negative; as the
pattern evolves, the feedback parameter becomes less negative, so the long-term sensitivity
exceeds the effective one.
\textbf{14.} $0.45 \times 44.009/12.011 = \qty{1.65}{\celsius}$ per \qty{1000}{PgC}.
\textbf{15.} 0.989 to \qty{2.31}{\celsius} per \qty{1000}{PgC}.
\textbf{16.} $2390 \times 12.011/44.009 = \qty{652}{PgC}$.
\textbf{17.} $0.45 \times 2.39 = \qty{1.08}{\celsius}$.
\textbf{18.} Within rounding, the carbon dioxide alone accounts for the whole human warming: the
warming by the other greenhouse gases and the cooling by the aerosols nearly cancel.
\textbf{19.} $1.5 - 1.07 = \qty{0.43}{\celsius}$.
\textbf{20.} $(0.43 - 0.1)/0.45 \times 1000 = \qty{733}{Gt}$ of carbon dioxide.
\textbf{21.} $733 - 210 = \qty{523}{Gt}$.
\textbf{22.} $26 \times 0.43 = \qty{11.2}{Gt}$; $523 - 11.2 = \qty{512}{Gt}$.
\textbf{23.} It falls by $0.1/0.45 \times 1000 = \qty{222}{Gt}$, to \qty{290}{Gt}: the
\qty{220}{Gt} by which the assessment says the choices about the other drivers move the budget.
\textbf{24.} Computed from the TCRE, the \omterm{def:b2:anthropogenic-climate-change:budget}{remaining carbon budget} for \qty{1.5}{\celsius} from the
start of 2020 is \textbf{\qty[text-series-to-math=true]{512}{Gt}} of carbon dioxide, within
\qty{3}{\percent} of the assessed \qty{500}{Gt} at a \qty{50}{\percent} probability.
\end{solution}
